\documentclass[11pt]{article}
\usepackage{amsmath,amssymb}
\title{Internal source transport: original tree words and their contacts}\author{}\date{}
\begin{document}\maketitle
\paragraph{Scope of the inputs.}
Equations (8)--(9) of the same coefficient schur cumulant decomposition calculation are restricted to its common-tree
fixed-log tensor. Scalar coefficient equality with FP coordinates
does not identify an arbitrary translated internal source.
In the conditional residual score contacts calculation, the two Gaussian Hermite contacts have the
correct subtraction and keep all Haar/quartic/cubic terms. Their
score variance is the residual classical energy after the actual
linear tree gauge map, not a substituted Euclidean norm.
Here is the missing source transport as a finite coefficient identity.
It is not a uniform nonlinear estimate or a bound on path lengths.

\paragraph{1. Exact source observable, with its cutoff.}
Let $q$ be the original canonical tree coordinates, with quadratic
precision $A$, original density $\rho_\alpha(q)$, and the cutoff
$\chi_\alpha(q)=\chi_{\rm phys}(\alpha q)$ from a fixed PHYSICAL
chart, equal to one near its flat minimum. Write $\chi=\chi_\alpha$
below. A fixed cutoff in canonical $q$ would have nonvanishing Gaussian
boundary mass and would NOT justify the full Gaussian coefficients.
For an INTERNAL direction
$w$, set $k=Aw$, $e=w^*Aw$. The change $q'=q+tw$ gives exactly
\[
 \begin{split}
 &\int\chi(q+tw)\rho_\alpha(q+tw)
             e^{tq^*k+t^2e/2}\,dq\\
 &\hspace{12mm}=\int\chi(q')\rho_\alpha(q')
             e^{tq'^*k-t^2e/2}\,dq'.
 \end{split} \tag{1}
\]
On the left, the Gaussian quadratic stays centered at $q$, while
the interaction, Haar factor AND cutoff are translated. Thus the
translated-vertex partition is the original observable on the right.
At each fixed finite block the cutoff-boundary terms are beyond all
Laplace coefficients. We do not discard them uniformly in block size.

In the local FP representation the observable is
\[
 \exp\{t k^*q_\alpha(z)-t^2e/2\}, \tag{2}
\]
where $q_\alpha(z)$ is the ACTUAL original tree quotient of
$U_b=\exp(\alpha z_b)$. It is not just its linearization.
The same cutoff is transported, and the FP density retains all gauge
noise, determinant and Haar factors of the local construction of the local gauge completion first nonlinear coefficient calculation.

\paragraph{2. A constructive original map.}
For the chosen original tree/forest, let $g_v(U)$ be its ordered
tree-path product from the retained anchor to $v$, with every anchor
frame fixed as in the original record. For each remaining link
$b=(s,t)$,
\[
 q_{\alpha,b}(z)=\alpha^{-1}\log
                  [g_s(U)U_b g_t(U)^{-1}]. \tag{3}
\]
Boundary roots and chord orientations are retained. Each bracket in
(3) is a literal ordered word $\prod_r\exp(\alpha X_r)$, with
$X_r=\pm z_{e_r}$. Define $q_\alpha=Lz+\alpha q_1(z)+\alpha^2q_2(z)+\cdots$.
The coefficients are fixed by the following explicit BCH recursion.
Starting with $s_1=s_2=s_3=0$, append a factor $X$ by
\[
 \begin{split}
 s_1^{\rm new}&=s_1+X,\\
 s_2^{\rm new}&=s_2+\tfrac12[s_1,X],\\
 s_3^{\rm new}&=s_3+\tfrac12[s_2,X]
   +\tfrac1{12}([s_1,[s_1,X]]+[X,[X,s_1]]).
 \end{split} \tag{4}
\]
Every right side uses OLD values. The resulting $s_1,s_2,s_3$ are
$Lz,q_1,q_2$. This gives actual frame-dependent source partners,
not an assumed abstract field redefinition. Their size may grow with
the original tree paths; no bound on that growth is asserted.

\paragraph{3. The transported covariance coefficient.}
Write the COMPLETE local FP density relative to its flat Gaussian as
$\exp[-\alpha Q_3-\alpha^2Q_4+O(\alpha^3)]$, where
$Q_3=T_3-\ell_1$ and $Q_4=T_4-J_2-\ell_2$.
Let $E$ denote its actual flat gauge-fixed Gaussian, and put
\[
 l_0=k^*Lz,\quad l_1=k^*q_1(z),\quad l_2=k^*q_2(z),\quad
 P=Q_4-\tfrac12Q_3^2,\quad m_1=El_1-E(l_0Q_3).
\]
At flat background $EQ_3=0$, and parity removes the order-$\alpha$
term of the covariance. Since $LGL^*=A^{-1}$, $El_0^2=e$.
Direct normalized expansion yields
\[
 \begin{split}
 [\alpha^2]\operatorname{Var}_{\rm FP}(k^*q_\alpha)
 ={}&E(l_1^2+2l_0l_2)-2E(l_0l_1Q_3)\\
 &-\operatorname{Cov}(l_0^2,P)-m_1^2 .
 \end{split} \tag{5}
\]
The normalization covariance and the mean square are essential.
By (1), the tree translated-vertex tensor is therefore
\[
 \boxed{K_{\rm tree}[w,w]=\operatorname{Cov}(l_0^2,P)
       +2E(l_0l_1Q_3)-E(l_1^2+2l_0l_2)+m_1^2.} \tag{6}
\]
This is the precise compensation for replacing the actual observable
$q_\alpha$ by its linearization. It includes ghost and Haar terms
inside $P,Q_3$ and the explicit original tree-coordinate contacts.

Set $v=GL^*Aw$. Gaussian integration by parts gives
$\operatorname{Cov}(l_0^2,P)=E D_v^2P$.
If the FP insertion is defined using variance rather than the square
of $Q_3$, then, at $EQ_3=0$,
\[
 K_{\rm FP}(v,v)=E D_v^2P+(E D_vQ_3)^2. \tag{7}
\]
Thus even the first term of (6) must be matched with its declared
mean convention. At flat SU(2), color symmetry can be used to verify
the zero tadpole before dropping this square; it is not silently
assumed in (6)--(7).

\paragraph{What this does not establish.}
This is a fixed-volume local source identity using the actual tree
words, original anchors, cutoff, and full FP density. It supplies no
uniform bounds for $q_1,q_2$ or for their signed combination in (6).
It does not identify a heat-cut marked tensor with the uncut density
expansion: cutoff regulator/source derivatives still have to be
transported. The conditional coarse compensation of the same coefficient schur cumulant decomposition calculation must also stay
attached. Full nonlinear background curvature and S94 remain open.
\end{document}
